% Imporditud: tools/import_eksperimendid.py
% Allikas: Eesti füüsikaolümpiaadi eksperimendi ülesannete
%   kogu (Rael Kalda, 2023), ülesanne Ü43, lahendus L43.
\setAuthor{EFO žürii}
\setRound{piirkonnavoor}
\setYear{2002}
\setNumber{G E1}
\setDifficulty{2}
\setTopic{dünaamika}
\setVahendid{joonlaud, lauatennise pall}
\setPunktid{10}
\setAllikas{Eksperimendi kogu Ü43, lahendus lk 53}
\setLahendusseis{toores}

\prob{Lauatennise pall}
Lauatennise pall kukub h = 80 cm kõrguselt kõvale horisontaalsele põrandale. Hinnata pallile mõjuva õhutaksitusjõu keskmist suurust, kui eeldada, et põrkel läheb 15\% palli kineetilisest energiast põranda ja palli siseenergiaks. Lauatennise palli mass on m = 2,5 g.

\hint

\solu
Idee --- mõõta kõrgus, milleni pall pärast põrget tõuseb (2 p). Puudu jäävale osale vastav potentsiaalne energia läks takistusjõu ületamisel tehtavaks tööks. Vastaja peab teadma kineetilise ja potentsiaalse energia valemeid ja töö valemit A = F s --- 3 p. Teeb kordusmõõtmisi ja leiab keskväärtuse --- 2 p. Avaldab õigesti otsitava suuruse --- 1 p. Hindab mõõteviga --- 2 p.
\probend
